\begin{resumenApertura}
El análisis concilia costos, precio y liquidez de la configuración 3.2, manteniendo el precio
aprobado para esta oferta sin reajuste de operación. Se separan la rentabilidad del proyecto,
la tesorería financiada y los beneficios del mandante.
\begin{recibe}
\item Evaluación reproducible, calendario de adquisiciones, estructura de soporte y valorización de innovaciones.
\end{recibe}
\end{resumenApertura}

\section{Costos frente a precio}\label{sec:eco-costo-precio}
La \tabref{eco-maestro} muestra los costos netos conciliados. Personal incluye los refuerzos de
producción y cierre de M21--M24; no se suman otra vez como incrementos. Los servicios externos
de montaje y campo se distinguen de la contratación interna. El hardware adquirido por el
mandante se informa separado de los egresos de audIT.

\tablaMontos*{Costo neto por fase}{eco-maestro}{costo-fase-implementacion,costo-fase-operacion,costo-proyecto}

La \tabref{eco-clasificacion} presenta el mismo costo neto por naturaleza para el Formulario T-22.

\tablaMontos*{Costo neto por naturaleza}{eco-clasificacion}{gasto-rrhh,gasto-inversion,gasto-operacion,gasto-administracion,gasto-total}

\montoPartida{gasto-proveedores} corresponde a infraestructura y servicios de operación. La
inversión incluye infraestructura y servicios de implementación.

La \tabref{eco-venta} presenta el precio por fase. El margen objetivo consolidado sobre ventas
es 18,50\,\%. La utilidad sobre costo del Resumen equivale a 22,70\,\%.

\tablaMontos{Precio neto ofertado por fase}{eco-venta}{implementacion,operacion}

La \tabref{hitos-e25} distribuye el precio de implementación en los doce hitos E-25.

\tablaHitosPago

La operación se factura de M21 a M56 y se cobra de M22 a M57 en el caso base.
La mensualidad neta es \montoPartida{cobro-m57}. La comisión de garantía
de funcionamiento se atribuye a implementación por objeto aun cuando se pague en M21.

La \tabref{eco-componentes} contrasta el costo y precio de cada componente. El margen se
calcula sobre ventas; los porcentajes de componentes no se suman para obtener el consolidado.

\begin{tabla}[ancho=completo,fuente=condensada]{Costo y precio por componente}{eco-componentes}{L{30mm} X X}
Componente & Costo neto & Precio neto \\ \midrule
Implementación directa & \montoPartida{cuadro-impl-directo-costo} & \montoPartida{cuadro-impl-directo-venta} \\
Implementación indirecta & \montoPartida{cuadro-impl-indirecto-costo} & \montoPartida{cuadro-impl-indirecto-venta} \\
Operación directa & \montoPartida{cuadro-op-directo-costo} & \montoPartida{cuadro-op-directo-venta} \\
Operación indirecta & \montoPartida{cuadro-op-indirecto-costo} & \montoPartida{cuadro-op-indirecto-venta} \\
\end{tabla}

La operación indirecta se traslada sin margen en esta clasificación. El margen directo compensa
esa distribución y mantiene 18,50 por ciento sobre la venta neta total. La plantilla RESUMEN
presenta además utilidad sobre costo: son bases de comparación distintas.

\section{Evaluación financiera}\label{sec:eco-evaluacion}
El VAN utiliza los flujos de la Fila A en M1--M56, con M0 igual a cero y tasa mensual 1,0261894\,\%.
La tasa pondera deuda al 0,9\,\% mensual y capital propio al 20\,\% efectivo anual supuesto.
La TIR mensual del caso base es 79,1905\,\%.

El VAN es \montoPartida{van}.
La cola de M57 se informa aparte. La recuperación descontada sostenida ocurre en M17.

La \tabref{eco-sensibilidad} conserva los escenarios calculados en la planilla. El déficit residual
mide faltante antes de cobros bajo el orden conservador de pagos; una caja mensual de cierre
positiva no demuestra suficiencia intramensual. Las TIR sin raíz convergente se rotulan para revisión.

\begin{tabla}[ancho=completo,fuente=condensada]{Sensibilidad y liquidez del escenario}{eco-sensibilidad}{X R{28mm} L{29mm} R{29mm}}
Escenario & VAN millones CLP & TIR mensual & Déficit millones CLP \\ \midrule
Atraso H5 M13 & 1.060,63 & Revisar raíces & 214,45 \\
Atraso H5 M14 & 1.056,65 & 20,18\,\% & 214,45 \\
H5 M14 + reserva M12 & 921,56 & 13,56\,\% & 368,65 \\
H5 M14 + reserva + facilidad & 921,56 & 13,56\,\% & 0,00 \\
Costos +5\% & 793,27 & 13,03\,\% & 53,21 \\
Cobros -5\% & 740,04 & 12,65\,\% & 24,06 \\
Reserva anticipada M9 & 917,36 & 11,86\,\% & 0,00 \\
H4 cobro M11, sin línea & 1.060,55 & 22,32\,\% & 245,25 \\
H4 cobro M11, con línea & 1.060,55 & 22,32\,\% & 195,25 \\
H4 cobro M11, facilidad M9 & 1.060,55 & 22,32\,\% & 0,00 \\
H4 acta y cobro M11 & 1.060,55 & 22,32\,\% & 160,18 \\
H12 cobro al mes siguiente & 1.060,99 & Revisar raíces & 0,00 \\
Soporte sin reducción & 964,81 & Revisar raíces & 0,00 \\
Soporte por madurez 3.0 & 1.064,66 & Revisar raíces & 0,00 \\
Amortización 3.1 en cuotas desde M21 & 1.064,66 & Revisar raíces & 0,00 \\
\end{tabla}

El equilibrio nominal es \montoPartida{equilibrio-nominal}.
Con el calendario base y escala uniforme, el umbral de venta a VAN cero es
\par\noindent\montoPartida{equilibrio-van}.

La facilidad contingente de \montoPartida{capital-adicional-h5} tiene desembolso previsto en M12
en el escenario de atraso H5 con reserva, sujeto a disponibilidad.
El capital inicial del caso base suma
\par\noindent\montoPartida{capital-inicial}.
Se financia en un 80\,\% con crédito bancario y un 20\,\% con capital propio. Los prepagos se
programan para M17 y M18. Las condiciones de financiamiento requieren respaldo comercial.

La reproducción del caso base se realizó en una copia de LibreOffice después de retirar los
valores guardados de las fórmulas. Los resultados se contrastaron contra el archivo 3.2. La
coincidencia del cálculo no valida por sí sola supuestos, contratos, capacidad ni una raíz
única de TIR. El criterio de decisión combina VAN, sensibilidad y liquidez.

\section{Curva S y flujo mensual}\label{sec:eco-curva}
La \figref{eco-curva-s} compara costos acumulados con intereses y facturación acumulada.
La \figref{eco-solvencia32} muestra por separado la caja financiada: es un saldo de tesorería,
no una tercera serie de avance contractual. La amortización del principal afecta la caja y no
se vuelve a contar como costo. Las dos figuras cubren M1--M56; M57 se concilia como cola de cobro.

\figura{economico/costos-venta/curva-s.pdf}{Costo con intereses y facturación acumulada}{eco-curva-s}
\figura{economico/costos-venta/solvencia-3.2.pdf}{Saldo de caja financiada del caso base}{eco-solvencia32}

La \tabref{eco-flujo-mensual} permite reproducir el flujo del proyecto y distinguirlo del
flujo financiado. Los valores se expresan en UF, con los parámetros convencionales E-24;
el modelo conserva las series exactas en CLP y sus conversiones. Fila A excluye financiamiento;
Fila C incorpora sus movimientos según las columnas del modelo. El saldo de caja es acumulado.

\begin{tabla}[ancho=completo,fuente=condensada]{Flujo mensual en UF}{eco-flujo-mensual}{L{9mm} R{26mm} R{26mm} R{26mm} R{26mm} X}
Mes & Cobro neto & Costo neto & Fila A & Fila C & Caja \\ \midrule
1 & 0,00 & 2.801,41 & -2.801,41 & 15.396,02 & 15.396,02 \\
2 & 8.955,13 & 1.931,90 & 7.023,22 & 8.593,30 & 23.989,32 \\
3 & 0,00 & 5.140,23 & -5.140,23 & -7.409,81 & 16.579,51 \\
4 & 7.835,74 & 3.662,21 & 4.173,53 & 5.392,22 & 21.971,72 \\
5 & 0,00 & 3.178,70 & -3.178,70 & -4.871,83 & 17.099,89 \\
6 & 11.193,91 & 4.827,04 & 6.366,87 & 7.541,05 & 24.640,94 \\
7 & 0,00 & 5.550,08 & -5.550,08 & -7.831,85 & 16.809,08 \\
8 & 0,00 & 5.381,54 & -5.381,54 & -5.537,33 & 11.271,76 \\
9 & 0,00 & 5.441,96 & -5.441,96 & -5.596,85 & 5.674,91 \\
10 & 11.193,91 & 4.734,11 & 6.459,80 & 8.448,06 & 14.122,96 \\
11 & 0,00 & 4.668,04 & -4.668,04 & -6.933,47 & 7.189,49 \\
12 & 11.193,91 & 4.633,04 & 6.560,87 & 8.549,12 & 15.738,61 \\
13 & 5.596,95 & 6.763,31 & -1.166,36 & -2.368,72 & 13.369,89 \\
14 & 5.596,95 & 4.937,34 & 659,62 & 520,68 & 13.890,57 \\
15 & 0,00 & 4.937,34 & -4.937,34 & -6.139,70 & 7.750,88 \\
16 & 11.193,91 & 5.132,64 & 6.061,27 & 8.046,67 & 15.797,54 \\
17 & 11.193,91 & 4.991,76 & 6.202,15 & -1.939,29 & 13.858,25 \\
18 & 11.193,91 & 4.403,32 & 6.790,59 & 121,15 & 13.979,40 \\
19 & 5.596,95 & 4.158,49 & 1.438,47 & 361,33 & 14.340,73 \\
20 & 0,00 & 3.787,98 & -3.787,98 & -4.865,12 & 9.475,61 \\
21 & 11.193,91 & 2.594,35 & 8.599,56 & 10.712,68 & 20.188,29 \\
22 & 2.701,94 & 2.232,80 & 469,14 & -1.671,42 & 18.516,87 \\
23 & 2.701,94 & 2.232,80 & 469,14 & 455,42 & 18.972,29 \\
24 & 2.701,94 & 2.260,68 & 441,25 & 392,08 & 19.364,38 \\
25 & 2.701,94 & 3.514,92 & -812,98 & -826,81 & 18.537,56 \\
26 & 2.701,94 & 2.084,82 & 617,12 & 603,29 & 19.140,85 \\
27 & 2.701,94 & 2.084,82 & 617,12 & 603,29 & 19.744,13 \\
28 & 2.701,94 & 2.090,19 & 611,74 & 597,91 & 20.342,04 \\
29 & 2.701,94 & 2.134,82 & 567,12 & 553,29 & 20.895,33 \\
30 & 2.701,94 & 2.187,32 & 514,62 & 500,79 & 21.396,11 \\
31 & 2.701,94 & 2.219,82 & 482,12 & 468,29 & 21.864,40 \\
32 & 2.701,94 & 2.084,82 & 617,12 & 603,29 & 22.467,69 \\
33 & 2.701,94 & 2.027,70 & 674,24 & 660,41 & 23.128,09 \\
34 & 2.701,94 & 2.027,70 & 674,24 & 660,41 & 23.788,50 \\
35 & 2.701,94 & 2.027,70 & 674,24 & 660,41 & 24.448,90 \\
36 & 2.701,94 & 2.173,13 & 528,80 & 483,81 & 24.932,71 \\
37 & 2.701,94 & 3.464,11 & -762,17 & -776,11 & 24.156,60 \\
38 & 2.701,94 & 2.034,01 & 667,92 & 653,99 & 24.810,60 \\
39 & 2.701,94 & 2.034,01 & 667,92 & 653,99 & 25.464,59 \\
40 & 2.701,94 & 2.039,39 & 662,55 & 648,62 & 26.113,21 \\
41 & 2.701,94 & 2.084,01 & 617,92 & 603,99 & 26.717,20 \\
42 & 2.701,94 & 2.136,51 & 565,42 & 551,49 & 27.268,69 \\
43 & 2.701,94 & 2.169,01 & 532,92 & 518,99 & 27.787,68 \\
44 & 2.701,94 & 2.034,01 & 667,92 & 653,99 & 28.441,68 \\
45 & 2.701,94 & 1.976,89 & 725,04 & 711,11 & 29.152,79 \\
46 & 2.701,94 & 1.976,89 & 725,04 & 711,11 & 29.863,90 \\
47 & 2.701,94 & 1.976,89 & 725,04 & 711,11 & 30.575,01 \\
48 & 2.701,94 & 2.003,17 & 698,77 & 653,77 & 31.228,78 \\
49 & 2.701,94 & 3.080,11 & -378,17 & -392,21 & 30.836,58 \\
50 & 2.701,94 & 1.983,21 & 718,73 & 704,70 & 31.541,28 \\
51 & 2.701,94 & 1.983,21 & 718,73 & 704,70 & 32.245,98 \\
52 & 2.701,94 & 1.988,58 & 713,36 & 699,32 & 32.945,30 \\
53 & 2.701,94 & 2.033,21 & 668,73 & 654,70 & 33.600,00 \\
54 & 2.701,94 & 2.085,71 & 616,23 & 602,20 & 34.202,20 \\
55 & 2.701,94 & 2.028,21 & 673,73 & 659,70 & 34.861,90 \\
56 & 2.701,94 & 2.353,21 & 348,73 & 334,70 & 35.196,60 \\
\end{tabla}


La cola de cobro de M57 se informa por separado y no se suma al VAN base M1--M56. Una caja
de cierre positiva no acredita disponibilidad diaria antes de los cobros; ese orden se revisa
en la sensibilidad de tesorería.

\section{Estrategia de adquisiciones}\label{sec:eco-estrategia}
La \tabref{eco-proveedores} relaciona las adquisiciones con su pagador y fecha de uso. Las referencias
de presupuesto no equivalen a una cotización aceptada ni a una recepción efectiva.

\begin{tabla}[ancho=completo]{Proveedores y condiciones de adquisición}{eco-proveedores}{L{32mm} L{21mm} L{25mm} X}
Proveedor o servicio & Pagador & Ventana & Condición de puesta en servicio \\ \midrule
Microsoft Azure por CSP & audIT & M1--M56 & Regiones Chile Central y Brazil South; circuito privado sujeto a factibilidad. \\
Dell y equipamiento de sala & audIT & M1--M6 & Energía, clima, extinción y recepción antes de H3. \\
Fortinet y hardware de campo & Curimón según T-11 & M5--M6 & Recepción y homologación del lote antes de montaje. \\
iWave, Teltonika y Technoton & Curimón según T-11 & Piloto M6, lotes M7--M9 & 182 kits iniciales, 210 balizas, cinco lectores; el stock de reposición no es montaje inicial. \\
Montaje por terminal & audIT & M7--M9 & 172 montajes externos; audIT conserva banco, recepción y prueba. \\
Ground Control / Iridium y LTE & audIT & Activación por unidad & Cobertura medida y activación solo al recibir el equipo. \\
GitLab y SonarSource & audIT & Según licencias & Cobertura comercial y condiciones de soporte confirmadas antes de usar. \\
Servicio de campo y soporte fabricante & audIT & Campo desde M19; soporte M13 & Visitas y rondas por orden; fabricantes con pagos anticipados anuales. \\
\end{tabla}

El proveedor masivo termina el montaje inicial en M9; los diez pilotos de M6 ya forman parte de
los 182 equipos. Si falta adhesión, stock o ventana se registra una desviación y se actualizan
instalación, activación y caja; no se da por ejecutado el escenario presupuestado.
La inversión referencial en hardware a cargo de Curimón suma \montoPartida{hardware-mandante}.
Las licencias a cargo de audIT totalizan \montoPartida{licencias-d3} en 56 meses.

Se adquiere equipamiento estándar y se conservan integración, banco, pruebas y recepción como
responsabilidad de audIT. La negociación busca disponibilidad, garantías, soporte y calendario
de pagos compatibles con H3 y los lotes. La dependencia de importación y homologación debe
confirmarse antes de emitir cada orden; no se afirma una internación ni una cotización recibida.

\section{Costos operacionales y perfiles}\label{sec:eco-operacion}

La operación combina personal dedicado, guardia, proveedores de infraestructura, nube, licencias,
conectividad, campo y administración. La \tabref{eco-operacion-categorias} muestra el costo por
categoría para los 36 meses. Las partidas corresponden a Cuadro maestro y RESUMEN, sin sumar
por segunda vez infraestructura ya incorporada a proveedores.

\tablaMontos*{Costo operacional por categoría}{eco-operacion-categorias}{final-op-1,final-op-2,final-op-3,final-op-4,final-op-5,final-op-6,final-op-7,final-op-8}


El precio de operación permanece sin reajuste. La reducción de carga variable por madurez no
se aplica a costos fijos ni autoriza rebajar cobertura SLA. La nómina y los servicios de
campo se controlan por separado.

\subsection{Costos y tarifas por perfil}\label{sec:eco-tarifas}

La \tabref{eco-tarifas-final} identifica costos y tarifas de referencia por fase, su perfil
y el rango de comparación E-26. Los roles adicionales se declaran como tales: la referencia
de rango no constituye una acreditación de competencias. Las tarifas desagregan el precio
cerrado de personal y no generan facturación horaria adicional al E-25.

\begin{tabla}[ancho=completo,fuente=condensada]{Costos y tarifas por fase en UF por hora}{eco-tarifas-final}{L{23mm} L{30mm} X R{24mm} R{24mm}}
Fase & Perfil & Referencia E-26 & Costo & Tarifa \\ \midrule
Implementación & Jefe de proyecto & Jefe de Proyecto & 1,45000000 & 1,75158405 \\
Implementación & Arquitectura & Arquitecto Senior & 1,10000000 & 1,32878790 \\
Implementación & Seguridad & Encargado de Seguridad TI & 1,10000000 & 1,50000000 \\
Implementación & Líder QA & Analista QA Experto & 0,60000000 & 0,80000000 \\
Implementación & Líder Datos & Ingeniero de Datos & 0,80000000 & 1,00000000 \\
Implementación & SRE & Ingeniero DevOps / SRE & 0,80000000 & 1,00000000 \\
Implementación & PMO & Analista / Diseñador Junior & 0,60000000 & 0,72479340 \\
Implementación & Liderazgo desarrollo & Jefe de Desarrollo & 1,05000000 & 1,26838845 \\
Implementación & Software & Ingeniero Senior & 0,75000000 & 0,90599175 \\
Implementación & Frontend & Ingeniero Senior & 0,75000000 & 0,90599175 \\
Implementación & Datos & Ingeniero de Datos & 0,80000000 & 1,00000000 \\
Implementación & QA & Analista QA Senior & 0,50000000 & 0,60399450 \\
Implementación & Firmware & Ingeniero Senior & 0,70000000 & 0,84559230 \\
Implementación & Electrónica & Ingeniero Senior & 0,70000000 & 0,84559230 \\
Implementación & Técnico de campo & Analista / Diseñador Junior & 0,40000000 & 0,50000000 \\
Implementación & Funcional & Analista / Diseñador Experto & 0,75000000 & 0,90599175 \\
Implementación & Integración & Ingeniero Senior & 0,75000000 & 0,90599175 \\
Implementación & Implantación y cambio & Analista / Diseñador Experto & 0,75000000 & 0,90599175 \\
Implementación & Analista terreno & Analista / Diseñador Junior & 0,50000000 & 0,60399450 \\
Operación & L1 24x7 & Analista / Diseñador Junior & 0,40000000 & 0,54678197 \\
Operación & SRE & Ingeniero DevOps / SRE & 0,80000000 & 1,09356394 \\
Operación & Seguridad & Encargado de Seguridad TI & 1,10000000 & 1,50365042 \\
Operación & Gestión servicio & Jefe de TIC & 1,05000000 & 1,43530268 \\
Operación & Monitoreo & TIC & 0,60000000 & 0,82017296 \\
Operación & Campo & Analista / Diseñador Junior & 0,40000000 & 0,54678197 \\
Operación & Software & Ingeniero Senior & 0,75000000 & 1,02521620 \\
Operación & Datos & Ingeniero de Datos & 0,80000000 & 1,09356394 \\
Operación & QA & Analista QA Senior & 0,50000000 & 0,68347746 \\
Operación & Arquitectura cierre & Arquitecto Senior & 1,10000000 & 1,50365042 \\
Operación & Frontend cierre & Ingeniero Senior & 0,75000000 & 1,02521620 \\
Operación & L1 suplencias y picos & Analista / Diseñador Junior & 0,40000000 & 0,54678197 \\
\end{tabla}


La suma de horas por tarifa conserva la venta de personal de RESUMEN en implementación y
operación. Se utilizan valores sin redondear en el cálculo; la presentación no modifica los
totales del contrato ni la base del impuesto.

\section{Mantención y soporte}\label{sec:eco-mantencion}

El presupuesto mantiene cobertura fija, guardia y la carga variable por tramo de madurez.
Las reducciones desde M25, M33 y M45 requieren carga observada compatible con SLA, cobertura
por turno y capacidad de escalamiento. No son una autorización para dejar puestos sin cubrir.
La \tabref{eco-soporte-final} identifica el costo horario y la dedicación inicial de los paquetes.

\begin{tabla}[ancho=completo,fuente=condensada]{Cobertura inicial de soporte}{eco-soporte-final}{X R{25mm} R{25mm} R{25mm}}
Paquete & Horas fijas & Horas variables & Costo UF/h \\ \midrule
L1 24x7 / EDT13.1 & 960,00 & 0,00 & 0,40 \\
SRE / EDT13.2 & 320,00 & 0,00 & 0,80 \\
Seguridad / EDT13.2 & 80,00 & 0,00 & 1,10 \\
Gestión servicio / EDT13.2 & 80,00 & 0,00 & 1,05 \\
Monitoreo / EDT13.2 & 160,00 & 0,00 & 0,60 \\
Campo / EDT13.3 & 80,00 & 0,00 & 0,40 \\
Software / EDT13.4 & 160,00 & 160,00 & 0,75 \\
Datos / EDT13.4 & 80,00 & 80,00 & 0,80 \\
QA / EDT13.4 & 80,00 & 80,00 & 0,50 \\
Arquitectura cierre / EDT6.5 & 0,00 & 0,00 & 1,10 \\
Frontend cierre / EDT6.5 & 0,00 & 0,00 & 0,75 \\
L1 suplencias y picos / EDT13.1 & 160,00 & 0,00 & 0,40 \\
\end{tabla}


Las horas de cierre M21--M24 se separan del régimen para evitar duplicarlas con mantenimiento
evolutivo. El campo externo por visita sustituye solo las tareas previstas y audIT conserva
coordinación, stock y recepción. La mantención preventiva se programa por activo, la correctiva
por incidente, y la evolutiva por bolsa priorizada con control de cambios y pruebas de regresión.
Los servicios de terceros y renovaciones se trazan a Proveedores, Detalle proveedores D4 y
Contrato de soporte. Los refuerzos por puesta en producción permanecen presupuestados.

\section{Valorización de innovaciones}\label{sec:eco-innovaciones}
La \tabref{eco-innovaciones} separa atribución de personal, egreso directo y beneficios de Curimón.
Las horas atribuidas pertenecen a la nómina existente y no constituyen un segundo gasto.

\begin{tabla}[ancho=completo]{Tratamiento presupuestario de las innovaciones}{eco-innovaciones}{L{39mm} X X}
Innovación & Imputación & Evidencia de resultado \\ \midrule
1. Expediente verificable & Firma \montoPartida{firma-inn1}; personal atribuido en implementación \montoPartida{innovacion-1-inversion} y operación \montoPartida{innovacion-1-costo}. & Medir expedientes y minutos evitados sin duplicar recuperos de INN-4. \\
2. Programación de patio & Desarrollo \montoPartida{innovacion-2-inversion}; operación \montoPartida{innovacion-2-costo} atribuidos. & Comparar preparación y primer intento contra línea base. \\
3. Acople y registro térmico & Desarrollo \montoPartida{innovacion-3-inversion}; operación \montoPartida{innovacion-3-costo} atribuidos. & El kilometraje es alcance base; medir acople y temperatura por separado. \\
4. Comodato y sobreestadías & 34 comodatos ya incluidos en 182 kits; nuevas altas requieren capacidad y pagador. & Recupero bruto de referencia \montoPartida{recupero-inn4}; no es cobro de audIT. \\
5. Prevención de fatiga & Personal atribuido en implementación \montoPartida{innovacion-5-inversion} y operación \montoPartida{innovacion-5-costo}; talleres y calibración \montoPartida{egreso-inn5} en Otros confirmados. & Medir resultado incremental y adopción; IVA y ejecutores a confirmar. \\
\end{tabla}

El caso base no incorpora beneficios hipotéticos de Curimón en los cobros de audIT. La atribución
mensual de cada innovación se verifica contra la capacidad disponible antes de reservar una tarea.

La \tabref{eco-innovaciones-montos} explicita inversión y operación atribuidas y el beneficio
reconocido en el caso base. La atribución de nómina no se agrega como un segundo gasto y el
beneficio del mandante no se transforma en venta de audIT.

\tablaMontos*[fuente=condensada]{Atribución económica de las cinco innovaciones}{eco-innovaciones-montos}{innovacion-1-inversion,innovacion-1-costo,innovacion-1-beneficio,innovacion-2-inversion,innovacion-2-costo,innovacion-2-beneficio,innovacion-3-inversion,innovacion-3-costo,innovacion-3-beneficio,innovacion-4-inversion,innovacion-4-costo,innovacion-4-beneficio,innovacion-5-inversion,innovacion-5-costo,innovacion-5-beneficio}

La medición incremental requiere línea base, adopción y responsable. Se mantiene beneficio
monetario reconocido igual a cero para audIT, mientras que los resultados de Curimón se
evalúan separadamente y sin doble conteo entre innovaciones.
